\scanpage{102}
Which also means that every \inserted{(or one)} splitting
of 31 --- for example defined by a point of $U_\eta$
\unsure{rational|rational} over $K$ --- induces a homomorphism $T\longrightarrow\pi$)

\noindent c) The action of $T$ on $\pi_{\mathrm{ab}}=H_1(U_{\bar\eta})$ is
trivial.

\noindent d) Likewise for $H_{\unsure{1|n}}(U_{\bar\eta},T_\ell)$, if $\ell$ is prime
to the residue characteristic. (\uline{NB} \deleted{\unsure{For characteristic $p>0$|For $\ell\ne p$|For $\ell=p$}} \inserted{$\ell=p>0$}\unsure{,|)}
one should be able to express it by a criterion
on the $p$-divisible group\ldots)

\noindent e) (In terms of a section of (30),
defining an extension of $\mathcal E_k$ by $\pi$,
and now assuming $k$ of finite
type over $\mathbf Q$) The said extension is
admissible (up to passing to an
open subgroup \inserted{of $\mathcal E_k$,} if need be), ie.
can be defined by an algebraic curve
over \inserted{$k$, or over} a finite extension of $k$.

\noindent f) (In terms of a section of (30)
and of one of (31)., allowing
$\mathcal E_k$ to act on $\pi$) : This
action is \inserted{not} \inserted{too} crazy (to be defined suitably !)% The end of convenablement continues on a short lower baseline; no hyphen is written.
\margin{101}
